\section{Deformation of the Whittaker counit}\label{sec:deformation}

Fix one of the groups $H$ under consideration.  Throughout this section and
Section~\ref{sec:plancherel}, put
\[
 Y=Y_H,\qquad b=\dim Y,\qquad D_s=D(L_H^s).
\]
Write $R_s$ for its enhanced Whittaker coefficient and $R'_s$ for the
opposite coefficient on $D_{-s}$.  We use the duality of
\eqref{eq:W2} to interpret $(R'_s)^\vee$, and define
\[
 J_s=F_H\circ(R'_s)^\vee[2b+2\delta_N].
\]
Thus $J_s$ applies a Whittaker functional to the first slot of the
$!$-diagonal kernel, whose two twists are $(-s,s)$.  Our objective is
\begin{equation}\label{eq:P1}
                         J_sR_s\simeq\operatorname{Id}_{D_s}
                         \qquad(s\notin\mathbb Q).
\end{equation}
The present section constructs the comparison and proves the finite
deformation statements needed to study it.  Section~\ref{sec:plancherel}
proves that its cone vanishes.

Since $F_H$ is an equivalence, it suffices to identify the kernels of
$J_sR_sF_H$ and $F_H$.  The latter is $\Delta_!k_Y$, with twists
$(-s,s)$.  Its first variable is a spectator, while $J_sR_s$ acts on the
second.  On a smooth affine chart $w:W\to Y$ in the first variable,
trivialize the determinant line and pull back the kernel.  Interchange
the factors and denote the resulting object on $Y\times W$ by $F_W$,
so the output variable now comes first.  Smooth base change describes
$F_W$ as the graph $!$-image, with
the smooth dimension conversion.  We construct the comparison on these
charts, naturally under smooth pullback, and then descend it to a map
of kernels.

\subsection{The ordinary-level inputs}\label{deform:classical-inputs}

We use the following established statements at exponent zero.  In each
statement the group is a connected reductive group over $\mathbb C$, the
curve is smooth and complete, the actual root datum is retained, and
the theta characteristic converts the half-density convention to the
ordinary $D$-module convention used here.

\begin{enumerate}
\item The geometric Langlands equivalence identifies $D(Y)$ with
$\operatorname{IndCoh}_{\operatorname{Nilp}}
(\operatorname{LS}_{\check H})$, where local systems are de Rham.
Under this identification the tempered inclusion is
\[
 \Xi:\operatorname{QCoh}(\operatorname{LS}_{\check H})
       \hookrightarrow
       \operatorname{IndCoh}_{\operatorname{Nilp}}
                         (\operatorname{LS}_{\check H}).
\]
The cone of the counit from Ran $!$-Poincar\'e applied to the enhanced
Whittaker coefficient lies in the kernel of the right adjoint to this
inclusion.  We call this kernel the \emph{anti-tempered category}.
The $!$-Poincar\'e images of the vacuum and its fixed-point Satake
translates are, up to one fixed shift and constant line, respectively
$\Xi(\mathcal O)$ and $\Xi$ of the tensor products of the evaluation
representation bundles.  These are the enhanced Whittaker and spectral
localization statements in the geometric Langlands theorem
\cite{GLCI,GLCII,GLCIII,GLCIV,GLCV}: Ran Casselman--Shalika identifies
$!$-Poincar\'e with spectral localization followed by $\Xi$, and the
right adjoint of spectral localization is fully faithful
\cite[Section~18.1.1, Proposition~18.1.6 and Equation~(18.4)]{GLCII}.
Together with \cite[Section~5.1 and Equation~(5.1)]{GLCI}, the
equivalence identifies the normalized counit with that of the
tempered inclusion.  In the
Drinfeld model the coefficient is $!$-pullback followed by Whittaker
projection; its left adjoint on holonomic finite-pole arguments is
$!$-image.

\item Ordinary geometric Satake and categorical Casselman--Shalika
identify the action of a Satake intersection complex on the vacuum
with the corresponding standard Whittaker object.  In the fusion
version these intersection complexes extend over the whole tuple
power of the curve, are universally locally acyclic over that power,
and specialize on diagonals to tensor products.  On a finite Hecke
correspondence they are relatively universally locally acyclic for
both projections; this follows locally from disk frames, with finite
jets sufficient in a fixed bound.  Pulling the correspondence to a
pole model gives the same action by changing the extension at the
marked points.  Both character signs are allowed: Cartan conjugation
over $\mathbb C$ interchanges them, also in families.

\item Lin's comparison between Poincar\'e series and miraculous
duality identifies the two geometric arrows in
\eqref{eq:P5} below, for ordinary holonomic Whittaker arguments.
We use the construction in
\cite[Propositions~2.2.4--2.2.5, Section~2.2.6,
Section~8.1.3 and Proposition~8.1.4]{Lin}, including
the version retaining the marked parameters.

\item The Whittaker characterization of tempering gives a finite
Whittaker test at a point that annihilates the anti-tempered category
\cite{FaergemanRaskin}.  Section~\ref{sec:plancherel} states the
precise finite-level functor and verifies its use in our situation.
\end{enumerate}

The passage from these ordinary-level inputs to irrational exponent uses
algebraically constructible calculations with a variable rank-one local
system.  Our deformation lemmas require mixed Hodge origin and
quasi-unipotent monodromy, so we first encode each Whittaker character by
an algebraic translation on an auxiliary affine line.  These translations
belong to the ordinary geometric diagrams to which the lemmas apply.
Fourier transform then turns them into the required exponential factors.

\subsection{An auxiliary axis and Fourier transform}

For a space $Z$ in a finite diagram, introduce $Z\times\mathbb A^1_u$.
We call a complex there a \emph{pre-complex}.  Fourier transform is
performed only after specialization of the coefficient parameter to a
field.  Our relative algebraic $D$-module Fourier transform is
\[
 \mathsf F(K)=p_{t,*}\bigl(p_u^!K\mathbin{\otimes^!}
                                      \operatorname{Exp}(ut)\bigr).
\]
Here the projections are from $Z\times\mathbb A^1_u\times\mathbb A^1_t$.
This normalization gives $\mathsf F(\delta_{u=0})=\omega_t$;
$\mathsf F[-1]$ is exact for the perverse convention.  The latter is
the Weyl-algebra Fourier theorem, applied relatively on affine charts
and with closed supports when the parameter scheme is singular.

\begin{lemma}[Fourier comparisons]\label{deform:fourier}
On holonomic arguments, the $!$ and $*$ integrals in the axis defining
$\mathsf F$ agree canonically.  Consequently Fourier commutes with
both pullbacks and both direct images for maps independent of the
axis, whenever these holonomic operations are defined.  These
identifications respect adjunctions, the $!$-to-$*$ map, and diagonal
orientation maps pulled back from the diagram without the axis.
\end{lemma}

\begin{proof}
Compactify the $u$-line and put $z=u^{-1}$ at infinity, retaining $t$
as a variable.  Extension across $z=0$ of the Fourier integrand is
clean.  Here is the local algebraic verification, which also explains
why retaining $t$ is necessary.  In a left-module presentation,
sections obtained from the module before tensoring with the
exponential, multiplied by Laurent functions of $z$ and polynomials
of $t$, are killed by powers of $z\partial_t-1$.  They generate the
meromorphic extension.  On a module supported at $z=0$, however,
$z\partial_t$ is locally nilpotent, so $z\partial_t-1$ is invertible.
The meromorphic star extension therefore has no quotient supported at
the boundary.  Duality excludes a boundary submodule of the
shriek extension.  Thus the cokernel and kernel, respectively,
of the clean comparison vanish.  Since both
extensions across the affine Cartier complement are perverse exact,
their comparison is an isomorphism.  The assertion for complexes
follows by perverse truncation.

Use the $*$ integral to commute $!$-pullback and $*$-image with
Fourier, and the canonically equal $!$ integral to commute $*$-pullback
and $!$-image.  Tensoring by the exponential connection is compatible
with either pullback.  Base change and the projection formula give
the resulting identifications.  Their transitivity, or equivalently
their adjunction descriptions, proves compatibility with units and
counits.  The same compactification identifies both constructions of
the $!$-to-$*$ map.  Finally a diagonal orientation class is external
to the axis; its pullback and tensor morphism use exactly these
adjunctions.  This proves compatibility also with orientation maps.
\end{proof}

We specialize Fourier transforms by $!$-pullback at $t=1$.  All
diagrams where this is done carry an algebraic torus action rescaling
$u$ through a nonzero power; the action may also move the other
variables.  Equivariance up to a Kummer factor suffices on frame
torsors.  Fourier gives the inverse action on $t$.  After a finite
cover of $\mathbb G_{m,t}$ and the coordinate change supplied by the
action, the diagram is an external product in $t$.  Base change at
$t=1$ therefore identifies its operations with the operations on the
specialized diagram.  For $*$ operations we use the smooth dimension
conversion in this product.  In particular,
$i_1^!\omega_t=k$, and covariant translation
$\operatorname{Trans}_{\psi}$ of the $u$-axis gives tensoring with
$\operatorname{Exp}(t\psi)$ after Fourier transform.  This is a
geometric base-change argument, rather than an interchange of a
nonproper image and an arbitrary specialization.

Let $p:Z\times\mathbb A^1_u\to Z$.  The projection that discards the
constant part in the axis is
\begin{equation}\label{eq:pre-axis-projection}
 C_!^u(K)=\operatorname{Fib}\bigl(K\longrightarrow p^!p_!K\bigr).
\end{equation}
The functor $p^!$ is fully faithful.  Its image Fourier-transforms to
objects supported at $t=0$, and the localization triangle gives
\[
 \mathsf F C_!^u(K)\simeq j_{t,!}j_t^!\mathsf F(K),
 \qquad j_t:\mathbb G_{m,t}\hookrightarrow\mathbb A^1_t.
\]
Thus $C_!^u(K)=0$ is the precise test for vanishing of Fourier on the
punctured axis.  A comparison whose cone is killed by $C_!^u$ will be
called an isomorphism \emph{modulo constants in $u$}.  Operations
independent of $u$, of either kind, preserve the constant subcategory,
as does translation by a function of the other variables.  Hence
these comparisons can be carried through subsequent operations.
We do not assert that $C_!^u$ commutes with a mixture of $!$ and $*$
operations.  When there are two axes the constructions and their
dilation actions are applied separately.

\subsection{Finite constructible calculations and Mellin ranks}

Put $A_0=\mathbb Q[q,q^{-1}]$, and let $M_q$ be the multiplicative
rank-one $A_0$-local system on $\mathbb G_m$ with monodromy $q$ and
its unit rigidification.  Twist calculations take place on frame
torsors.  A factor $M_q(r)$ means ordinary pullback and ordinary
tensor product; equivalently one can use its dualizing normalization
with $!$-tensor product.  Specialization at
$q=\exp(2\pi i s)$ and Riemann--Hilbert, after complexification,
recover the regular-holonomic calculation with exponent $s$.

We make explicit the finiteness scope used below.  A finite geometric
calculation consists of a finite diagram of bounded algebraically
constructible complexes, built by the six operations, cones, and
finite totalizations from constant complexes, intersection complexes,
and the displayed Kummer factors.  The maps in the diagram are the
geometric maps of these operations.  Over an output chart the maps
are finite type; maps used for $!$-image are schematic or separated
representable.  Smooth unipotent-extension stack projections are
also permitted.  Their two-term vector-stack presentations reduce
the calculation to schematic ones, with the specified shifts and
Tate twists.

For $*$-image we also allow finite-type Artin stacks with unipotent
identity components of stabilizers and finite component groups.
Stratify such a stack into gerbes over algebraic spaces, after
flattening inertia and rigidification.  Smoothly or \'etale locally
these gerbes have the calculation of a classifying stack.  A
connected unipotent group is constructibly acyclic, and pullback
along its torsors is fully faithful; one may further stratify to
use split unipotent coordinates.  Finite component groups contribute
exact invariants in characteristic zero.  Gluing the strata by
localization triangles proves bounded constructibility and the
coefficient formulas.  The same argument, using safe base change,
applies to holonomic $D$-module $*$-images.  It also extends
Lemma~\ref{deform:fourier} to the stack operations just listed.
Locally finite-type non-quasi-compact stacks are used only through
these relative finite constructions over charts.

Every stalk of a finite calculation is a perfect $A_0$-complex, and
coefficient specialization is derived tensor product.  For schemes
these facts follow by finite algebraic triangulations and the same
finite link calculations for pullback and extension.  The reductions
above establish them for the allowed stacks.

We shall also use the mixed enhancement of these diagrams.  Saito's
theory supplies the six operations, pure intersection complexes,
weight inequalities, and quasi-unipotence of nearby cycles
\cite{SaitoMHM}; the enhanced formalism, including stacks, supplies
the coherent diagrams of operations and their comparisons
\cite[Sections~1 and~3]{TubachMHM}.  All weight conclusions will be
tested on finite schematic charts.  In particular, cones and
adjunction maps below are formed in the mixed category, not by
choosing arbitrary maps after forgetting the mixed structure.
Our weight convention is
\[
 K\text{ has weight }\le w
 \quad\Longleftrightarrow\quad
 {}^pH^j(K)\text{ has weights }\le w+j\text{ for every }j.
\]
The dualizing complex of a smooth scheme has weight zero.
Ordinary pullback and $!$-image preserve upper weight bounds, as does
the smooth conversion $[2e](e)$ in relative dimension $e$.

\begin{lemma}[Non-torsion rank test]\label{deform:rank}
Consider a finite geometric diagram as above.  Assume that its
$q$-dependence is introduced by Kummer factors $M_q(r)$ for algebraic
invertible functions $r$, and that the corresponding diagrams with
these factors recorded as graphs have mixed Hodge origin.  At every
output stalk, the dimensions of its specialized cohomologies and
the ranks of the induced maps are constant as $q$ ranges over
complex numbers that are not roots of unity.  In particular, a
finite output complex which vanishes at $q=1$ vanishes at every such
$q$.
\end{lemma}

\begin{proof}
To compare non-torsion specializations, introduce a recording coordinate
$z\in\mathbb G_m$, independent of the Fourier axis.  Replace an
initial factor $M_q(r)$ by the graph
$z=r$ and apply at the end
\begin{equation}\label{eq:P2}
                K\longmapsto p_!\bigl(K\otimes M_q(z)\bigr),
\end{equation}
where $K$ is independent of $q$ and $p$ forgets $z$.
For a complex of mixed Hodge origin, the nearby-cycle monodromies
at $z=0,\infty$ are quasi-unipotent.  Tensoring with $M_q(z)$ for
non-torsion $q$ leaves no eigenvalue $1$ there.  Thus both invariants
and coinvariants vanish, and the extension over these endpoints is
clean.  Consequently the $!$ integral in \eqref{eq:P2} is canonically
the $*$ integral.  This criterion holds relatively, with charts and
spectators retained.  The argument of Lemma~\ref{deform:fourier}
therefore moves \eqref{eq:P2} past either pullback, either direct
image, and their geometric comparison maps.

For external products multiply the recording coordinates and use
$!$-image along multiplication, since $M_q$ is multiplicative.  This
also explains cancellation of ratio twists.  Before
\eqref{eq:P2}, a frame change rescales the recording coordinate by
the corresponding positive or negative power.  Two occurrences of
an inner frame have opposite scalings.  Their product recording
coordinate is unchanged, so the complex descends through that
frame torsor before the inner frame is removed.  In local
trivializations this is simply invariance of the product under
opposite changes of variables by a unit.  The graph construction
and its maps respect these equivariances.  We have therefore obtained
\eqref{eq:P2} as a presentation of the whole specialized diagram for
non-torsion $q$, including its maps.  At $q=1$ and with logarithm
coefficients, we retain the original geometric diagram: the recording
presentation is used only on the non-torsion locus.

Pull to an output stalk.  If $P$ is perverse on $\mathbb G_m$ and
of mixed Hodge origin, Artin vanishing and equality of $!$ and $*$
integration imply that
$R\Gamma(\mathbb G_m,P\otimes M_q)$ is concentrated in degree zero.
Its dimension is its Euler characteristic, independent of $q$:
stratify $P$ and use that tensoring with a rank-one local system
does not alter the Euler characteristic.  The same reasoning
applies to mixed subquotients.  Taking perverse cohomology now
proves constancy in each degree for a complex and exactness on the
perverse subquotients.  It also proves constancy for maps, either
by their images in perverse cohomology or by the long exact
sequence of their cones.  Finally, perfectness of the original
$A_0$-complex and vanishing of its derived fibre at $q=1$ give
vanishing at the generic parameter.
The constant-rank conclusion on the non-torsion locus gives the
last assertion.
\end{proof}

The preceding lemma concerns finite calculations.  A zero colimit at
$q=1$ says that each cohomology class dies under some later transition;
it does not say that any finite-stage complex vanishes.  We therefore
cannot apply the finite rank test directly to an increasing collection
of pole bounds.  The next lemma uses a uniform upper weight bound to
turn vanishing at a finite logarithm length into a transition that
kills the entire completed cohomology module.

\subsection{Logarithm coefficients and passage to colimits}

Write $B_n=\mathbb Q[h]/h^n$.  Substitution $q=\exp(h)$ identifies
$A_0/(q-1)^n$ with $B_n$.  The length-$n$ logarithm variation
$\mathcal L_n$ on $\mathbb G_m$ has this monodromy and the usual
Tate convention for $h$.  Its quotients and ideals fit into exact
sequences of admissible variations
\begin{equation}\label{eq:logarithm-sequence}
 0\longrightarrow\mathcal L_{n-a}(a)
  \xrightarrow{\ h^a\ }\mathcal L_n
  \longrightarrow\mathcal L_a\longrightarrow0
                         \qquad(n\ge a),
\end{equation}
with the first term zero when $n=a$.  One can construct these
variations as symmetric powers of the Kummer extension; their
successive quotients are $\mathbb Q(j)$, $0\le j<n$.
In particular multiplication by $h$ lowers weights by two.

\begin{lemma}[Logarithm colimit test]\label{deform:colimit}
Let $\{K_i\}$ be a filtered system of finite geometric calculations
with outputs on a fixed finite schematic chart.  Suppose that each
calculation and all transition maps are linear in a single input
$M_q(r)$.  Assume their substitutions by $\mathcal L_n(r)$ lift to
mixed diagrams compatibly with \eqref{eq:logarithm-sequence}.
For every stalk $x$ and every degree $d$, suppose that
\[
 \underset{i}{\operatorname{colim}}\,
 H^d\bigl((K_i)_{x,q=1}\bigr)=0
\]
and that the weights of these groups have a common upper bound
$C_{x,d}$, independent of $i$.  Then the same cohomology colimit
vanishes at every non-root-of-unity value of $q$.
\end{lemma}

\begin{proof}
Fix $x,d$ and write $C=C_{x,d}$.  Denote the stalk calculation with
length-$n$ logarithm coefficients by $C_{i,n}$.
Linearity makes \eqref{eq:logarithm-sequence} an exact triangle of
these calculations.  Induction on $n$, the long exact sequence,
and exactness of filtered colimits show that
\[
 \operatorname*{colim}_i H^j(C_{i,n})=0
                  \quad\text{for every }j,n.
\]
The same filtration shows that the weights of $H^d(C_{i,n})$ are
at most $C$ for all $i,n$: each successive term comes from
$H^d(C_{i,1})(a)$ and has weights at most $C-2a$.

For a fixed $i$, perfectness and the Mittag--Leffler property give
\[
 M_i:=\varprojlim_n H^d(C_{i,n})_{\mathbb C}
   =H^d\bigl((K_i)_x\mathbin{\otimes}^{L}_{A_0}
                                             \mathbb C[[h]]\bigr).
\]
This is a finite $\mathbb C[[h]]$-module.  The coefficient exact
sequence identifies $M_i/hM_i$ with the image of $M_i$ in
$H^d(C_{i,1})_{\mathbb C}$; it need not be all of that group.
For precision, these images carry canonical mixed Hodge
structures.  Put
$P_i=(K_i)_x\otimes^L_{A_0}\mathbb C[[h]]$, and choose $e$
such that $h^e$ annihilates the $h$-power torsion of
$H^{d+1}(P_i)$.  The universal
coefficient sequence is
\[
 0\longrightarrow M_i/h^nM_i\longrightarrow
 H^d(C_{i,n})_{\mathbb C}\longrightarrow
 H^{d+1}(P_i)[h^n]\longrightarrow0.
\]
Under transition from length $m$ to length $n$, the last map
on these torsion terms is multiplication by $h^{m-n}$.
For $m\ge n+e$ the image is therefore exactly $M_i/h^nM_i$:
containment follows from that vanishing, and surjectivity
onto it follows from $M_i/h^mM_i\to M_i/h^nM_i$.
Thus the latter modules are stabilized-image subobjects of
the finite-length mixed Hodge structures, with surjective
mixed transition maps.

Deligne's functorial complex splitting of mixed Hodge structures
\cite{DeligneHodge} permits a finite set of generators of this
image to be chosen homogeneous for the weight splitting and
lifted to weight-homogeneous elements of the inverse system.
Indeed, choose a basis in Deligne's bigrading at length one
and lift recursively in the same bigraded components of
the stabilized images.  Strictness gives surjectivity on
each such component.  The resulting elements lift
coherently and remain homogeneous in weight.  Nakayama's
lemma says that their pro lifts generate $M_i$.

Let $v$ be such a pro generator, of weight $w$.  Choose $a$ with
$C-2a<w$.  Since the length-$a$ colimit is zero, a transition
$i\to i'$ kills the length-$a$ component of $v$.  For every
$n\ge a$, exactness of \eqref{eq:logarithm-sequence} puts its
image at length $n$ in the image of
\[
                 H^d(C_{i',n-a})(a)
                       \longrightarrow H^d(C_{i',n}).
\]
The source has weights at most $C-2a$, whereas the image of $v$
is homogeneous of weight $w$.  Strictness forces that image to
vanish.  Its components at lengths below $a$ vanish as well.
Thus the whole pro generator is killed by this transition.
There are finitely many generators, so one further transition
kills $M_i$.

The substitution $q=\exp(h)$ embeds $\mathbb C(q)$ into
$\mathbb C((h))$: a nonzero Laurent polynomial cannot
vanish as a formal exponential polynomial.  After extension
to $\mathbb C((h))$, the preceding conclusion says that the
cohomology map from stage $i$ to a sufficiently late stage is
zero at the generic coefficient parameter.  Lemma~\ref{deform:rank}
applied to this finite transition makes its rank zero at every
non-torsion $q$.  Each stage in each degree is consequently
killed later in the system, which is precisely the required
vanishing of the filtered colimit.
\end{proof}

\begin{lemma}[Stalk detection for the colimits used here]
\label{deform:stalks}
On a finite smooth chart, the filtered colimits of regular-holonomic
complexes occurring above are detected by constructible stalks.
The analogous holonomic Ind-category has a separated perverse
$t$-structure and embeds fully faithfully into all $D$-modules.
\end{lemma}

\begin{proof}
Bounded holonomic complexes on the chart are compact as
$D$-modules.  Hence their Ind-category, and likewise the Ind-category
of regular-holonomic complexes, embeds fully faithfully.  The
perverse $t$-structure extends to a separated $t$-structure there.
For regular-holonomic complexes, the bounded comparison between
the perverse and ordinary constructible $t$-structures passes to
Ind, so the latter is separated as well.

It remains to justify stalk detection in its heart.  The ordinary
algebraically constructible heart is Noetherian.  Indeed, consider
an ascending sequence of subsheaves of one constructible sheaf.
Choose a smooth connected open stratum where the ambient sheaf
is lisse.  The generic ranks stabilize.  A term of that rank is
lisse as a subsystem on a smaller open.  Every later term agrees
with it there: the quotient is a subsheaf of a local system with
zero generic rank, hence is zero.  Induction on the dimension of
the closed complement proves stabilization.  A finite
stratification deals with several top-dimensional components.
Every object of the Ind-heart is a filtered union of its
Noetherian subobjects.  If all its stalks vanish, each such
subobject has zero stalks and hence is zero.  Stalk detection
and separatedness prove the claim for complexes.  We have used
constructible stalks only inside this Ind-category, not as a
test for arbitrary all-$D$-module objects.
\end{proof}

\subsection{Finite Whittaker stages}\label{deform:finite-stages}

We now apply the deformation lemmas to the kernel $F_W$.  Let
$S_I=X^I$, where $I$ is a nonempty finite set, and choose a common
pole bound $m$ at its marks.  Write
\[
 \pi:\mathcal Y_m\longrightarrow Y,\qquad
 \pi_S:\mathcal Y_m\longrightarrow S_I
\]
for the pole model of Section~\ref{sec:whittaker}; when appropriate
$\pi$ retains the marked positions.  Spectator variables are
carried throughout.  For a relevant stratum $\sigma$, write
$f_\sigma,q_\sigma,\psi_\sigma$ for its embedding, unipotent-stack
projection, and character.  The base of $q_\sigma$ is the
corresponding configuration open $S_\sigma$ in a partition
stratum of $S_I$, and set $e_\sigma=\dim(q_\sigma)$.

Lift $F_W$ to its constructible graph formula on frame torsors.
Its only coefficient input is $M_q(r)$, for the relative frame
ratio determined by the chosen spectator trivialization.  The
formula is the smooth pullback of $\Delta_!k_Y$, so its smooth
dimension shifts have the corresponding mixed Tate twists.
Put
\[
             \widehat F=\pi^!(F_W\boxtimes\delta_{u=0}).
\]
For the sign giving coefficient $\operatorname{Exp}(\psi)$, define
\begin{equation}\label{eq:P3}
 \begin{split}
 \widehat T_\sigma&=f_{\sigma,!}
             \operatorname{Trans}_{+\psi_\sigma}q_\sigma^!,\\
 \widehat H_\sigma&=q_{\sigma,*}[-2e_\sigma]
             \operatorname{Trans}_{-\psi_\sigma}f_\sigma^!.
 \end{split}
\end{equation}
In mixed calculations the shift in the second line includes
$(-e_\sigma)$.  Filtered determinants identify the twists along
these operations.  They therefore lift to the same frame
torsors and add no new Kummer input.  The opposite sign replaces
$\psi$ by $-\psi$ everywhere.

\begin{proposition}[Finite projection diagrams]\label{deform:stages}
There is a canonical finite triangle
\[
 \widehat M_{I,m}\longrightarrow\widehat F
                  \longrightarrow\widehat F_\perp
\]
such that $\widehat M_{I,m}$ is a finite extension of images of the
$\widehat T_\sigma$ and every $\widehat H_\sigma$ kills
$\widehat F_\perp$.  After field specialization, Fourier transform,
and $!$-restriction at $t=1$, its first map is the counit of
Whittaker projection of $\pi^!F_W$.  Denote that projection by
$M_{I,m}$.  The augmented stages
\begin{equation}\label{eq:P4}
        \pi_!\widehat M_{I,m}
                    \longrightarrow F_W\boxtimes\delta_{u=0}
\end{equation}
form a geometric diagram under increases of pole bounds and
closed diagonals of marked parameters.  This entire diagram is
linear in the single initial Kummer input.
\end{proposition}

\begin{proof}
The stratum calculation of
Proposition~\ref{whittaker:strata-properties} gives the adjunction
$\widehat T_\sigma\dashv\widehat H_\sigma$ and
$\widehat H_\sigma\widehat T_\sigma=\operatorname{Id}$.
There are only finitely many relevant strata in a pole bound.
Order them from the higher open strata to the lower ones.
For each successive stratum, take the cone of its projection
counit on the current remainder.  Lower shriek extensions have
zero restriction to the higher open strata already treated.
Induction therefore kills all the $\widehat H_\sigma$ and
produces the stated triangle.  Orthogonality to the generated
subcategory characterizes the triangle, so its construction is
independent of the chosen refinement of the order and is
functorial, with coherent comparison maps.

Cartan scaling by $2\rho$ acts on the Pl\"ucker data and scales
each $\psi_\sigma$ by the square (or inverse square in the other
action convention).  Give $u$ the same scaling.  The diagram is
equivariant for this action.  On the twist line pulled from $Y$
the action is trivial; its identification with a stratum line
over $S_\sigma$ carries the induced filtered-determinant action.
Thus the required dilation equivariance persists in frame
presentations containing Kummer factors.  The Fourier
specialization rules apply.

After specialization, $\widehat T_\sigma$ becomes the corresponding
Whittaker shriek extension.  It satisfies character equivariance
on the stratum, and its holonomic extension preserves that
equivariance by the smooth-arrow tests in \eqref{eq:W1}.
Vanishing of all the specialized $\widehat H_\sigma$ is precisely
the stratum criterion for zero Whittaker projection.
Consequently the triangle specializes to the claimed projection
triangle.  The object $M_{I,m}$ is locally holonomic.  It is
supported cleanly on a bounded-zero quasi-compact open containing
the relevant strata: its $!$-restriction to the complement is
zero by the stratum criterion; its $*$-restriction is zero by
holonomic duality and the same criterion for the opposite
character.  This support conclusion is asserted after Fourier
specialization; the pre-complex need not have both kinds of
vanishing support.

Apply $\pi_!$ to the first arrow and use its adjunction with
$\pi^!$ to obtain \eqref{eq:P4}.  These are finite geometric
operations on every output chart.  The map $\pi$ in a pole
bound is relatively finite type; alternatively the first term
is computed by its finite extensions from quasi-compact
relevant strata.

For transitions use the closed inclusions that increase pole
bounds and the closed diagonals $S_I\hookrightarrow S_J$
associated to surjections $J\twoheadrightarrow I$.  Increase
the new pole bound enough to include sums at coincident marks.
The $\widehat F$'s restrict by $!$ to each other.  Closed image
preserves the subcategory generated by the
$\widehat T_\sigma$: for pole inclusions the old relevant strata
are strata at the new bound, and for diagonals coincidences
are already part of the stratification.  The earlier
$\widehat M$ therefore maps to the later $\widehat F$ through
the later $\widehat M$.  Orthogonality gives the factorization
and all its coherences.  Every step is geometric and linear
in the single $M_q(r)$; hence the logarithm lifts and the
recording presentation \eqref{eq:P2} apply to the diagram.
\end{proof}

For clarity, its colimit can be calculated without a compactness
claim about the kernel.  Let $r_{I,m}$ be the $!$-restriction
functors in the limit presentation of the Whittaker category,
and $a_{I,m}$ their left adjoints in the corresponding colimit
presentation.  For an object $A$ in that category,
\[
                A\simeq\operatorname*{colim}_{I,m}
                                  a_{I,m}r_{I,m}A.
\]
Indeed mapping this formula to $B$ gives, by adjunction, the
limit of the mapping complexes between $r_{I,m}A$ and
$r_{I,m}B$, which is the mapping complex in the limit category.
Here $M_{I,m}$ computes $r_{I,m}R_sF_W$.  At $s=0$, applying
the known left adjoint $!$-Poincar\'e to
$a_{I,m}M_{I,m}$ gives $\pi_!M_{I,m}$, again by adjunction.
Both classical functors are continuous, so their tensored
adjunction permits the spectator; the partial $!$-images on
the holonomic arguments are their actual left adjoints there.
Thus the Fourier-specialized colimit of \eqref{eq:P4} at
exponent zero is exactly the ordinary $!$-Poincar\'e counit.

We will apply Lemma~\ref{deform:colimit} to finite calculations
presenting this colimit.  Take the simplicial replacement on a
skeleton of the stage diagram, restrict to finite subdiagrams,
and then to finite skeleta.  These form a filtered system with
the same colimit.  In the skeletal filtration, simplicial
degree $l$ contributes with shift $[l]$.  Taking the augmented
cone gives a filtered system $\{\widehat C_\beta\}$ of finite diagrams,
where $\beta$ records the finite subdiagram and skeletal truncation.
These cones have all the coefficient and mixed-lift properties just
proved.  The shift $[l]$ will be included in the weight estimate in
Section~\ref{sec:plancherel}.

\subsection{Comparison through the compactified diagonal}

We next identify each finite Fourier-specialized stage with
the formula defining $J_s$.  This is a finite comparison, so
Lemma~\ref{deform:rank}, rather than the colimit lemma, suffices.
Recall the proper Vinberg correspondence from
Section~\ref{sec:pseudoidentity},
\[
 b_{\mathrm V}:\bar Y\longrightarrow Y\times Y,
 \qquad J^b=(j_0)_!l_*k_Y,
\]
where $J^b$ carries the ratio twist with exponent $-s$ in the
first slot.  For a locally holonomic Whittaker object $D$ in
a pole bound, form
\[
 \mathscr Z=\mathcal Y_m\times_Y\bar Y,
 \qquad g:\mathscr Z\longrightarrow\mathcal Y_m\times\bar Y,
 \qquad p_2:\mathscr Z\longrightarrow Y.
\]
The last map uses the second factor of the Vinberg
correspondence.  All these formulas retain spectator
variables, and initially also the marked parameter $S_I$.

Let $\mathscr Z^{\mathrm{tr}}$ be the open where the Borel of
the Pl\"ucker data and the opposite parabolic of the Vinberg
data are transverse at the generic point of $X$.  It contains
the data over $j_0$.  It is open because, in the section-matrix
description, transversality is detected by generic
nonvanishing of the Vinberg maps on the Pl\"ucker lines.
The open-orbit description of each Vinberg fibre gives the
test on fundamental lines for the omitted roots; for the
actual root datum take positive multiples of the relevant
weights.  Transversality then implies nonvanishing for every
highest-line map.

\begin{lemma}\label{deform:transverse-map}
The restriction $p_{2,\mathrm{tr}}$ is representable and
separated.  On finite output charts the unrestricted
$p_{2,*}$ has the unipotent-stabilizer finiteness required for
the finite geometric calculus.
\end{lemma}

\begin{proof}
At the generic point, an automorphism of the Pl\"ucker data
lies in the maximal unipotent.  Nonzero Vinberg matrices on
the highest lines force every residual torus scaling to be
trivial.  The remaining kernel of the Vinberg map is in the
radical of the opposite parabolic, whose intersection with
this maximal unipotent is trivial.  Automorphisms over the
output are therefore trivial, proving representability.

For separatedness use the valuative criterion for isomorphisms.
The torus scalings comparing the matrices extend as units:
their compositions on the Pl\"ucker lines are generically
nonzero on the special curve.  At the valuation given by the
generic point of that special curve, choose integral frames
putting both systems of Pl\"ucker vectors in standard form.
This is possible because the reductions are nondegenerate.
The isomorphism over the fraction field is now an element
$n$ of the maximal unipotent.  Choose an actual strictly
dominant highest-weight representation, highest vector $e$,
and an output covector whose composition with the Vinberg
matrix takes unit value on $e$.  The resulting row $\ell$
is integral, $\ell(e)$ is a unit, and $\ell n$ is integral.

These conditions force $n$ to be integral.  To verify this,
write its ordered root coordinates by increasing positive
root height.  Evaluate $\ell n$ on $f_\beta e$, where
$f_\beta$ lowers by a positive root $\beta$.  The coordinate
for $\beta$ occurs with coefficient a nonzero root-string
integer times $\ell(e)$, hence a unit over the complex DVR.
All other nonconstant terms involve coordinates of smaller
height: the total raising height cannot exceed that of
$\beta$, and a different root of the same height cannot
reach the highest weight.  Induction recovers each
coordinate integrally.  The isomorphism consequently
extends at the generic point of the special curve.
Normality and the codimension-one extension criterion for
sections of the affine isomorphism scheme extend it over
the entire curve family; its compatibility with the data
extends as well.  The argument retains marked parameters;
if they are forgotten, their own separatedness supplies
that part of the criterion.

Finally, before restricting to the transverse open, inertia
over an output chart embeds in the inertia of the
Pl\"ucker data, because the Vinberg compactification is
representable.  The latter is unipotent.  The previously
established unipotent-stabilizer reduction therefore
applies to $p_{2,*}$.
\end{proof}

Consider the two geometric arrows
\begin{equation}\label{eq:P5}
\begin{split}
 p_{2,\mathrm{tr},!}
       \bigl(g^*(D\boxtimes J^b)|_{\mathrm{tr}}\bigr)
 &\longrightarrow
 p_{2,\mathrm{tr},*}
       \bigl(g^!(D\boxtimes J^b)|_{\mathrm{tr}}\bigr)[2b]
 \\[-2pt]
 &\longleftarrow
 p_{2,*}g^!(D\boxtimes J^b)[2b].
\end{split}
\end{equation}
The first arrow uses the orientation of the smooth diagonal
of $Y$ and then the $!$-to-$*$ comparison on the transverse
open.  More explicitly, pull the diagonal class to the
fibre product by the canonical base-change map and use the
tensor-to-$!$-pullback morphism; this gives
$g^*\to g^![2b]$.  The tensor morphism also follows from
adjunction and the projection formula.  The second arrow
is restriction to the transverse open.  Mixed versions
include the Tate conversion belonging to $[2b]$.

These are finite calculations on an output chart.  Indeed,
cut $\mathcal Y_m$ to a bounded-zero quasi-compact open
containing the relevant strata.  The clean support of $D$
just proved shows that the star pull in the first term
extends by shriek from this cut, whereas the shriek pulls
in the other terms extend by star.  Thus all three images
are computed from the cut.  Properness of the Vinberg
correspondence now gives the finite calculation above an
output chart.  Enlarging the cut gives the same arrows by
restriction and transitivity of the $!$-to-$*$ comparison.
Proper base change and compatibility of orientation also
make them natural under closed pushes in the stage indices.

After integrating the proper marked parameter $S_I$, the
left term of \eqref{eq:P5} is $\pi_!D$: this is the
shriek calculation of the diagonal using the Vinberg
correspondence.  The right-source term is $J_s(a_{I,m}D)$.
For the latter identification, \eqref{eq:W2} transposes
$a_{I,m}$ to restriction.  Pairing with $D$ makes the
opposite Whittaker-projection counit an isomorphism, so
one can use the unprojected pullback of the $!$-diagonal
kernel.  Vinberg proper pushforward and $!$-base change
then give exactly the right-source expression.  The
$2\delta_N$ in $J_s$ cancels the normalization in
\eqref{eq:W2}; the remaining diagonal shift is $2b$.

\begin{proposition}[Finite Vinberg comparison]\label{deform:vinberg}
For $s\notin\mathbb Q$ and $D=M_{I,m}$, both arrows in
\eqref{eq:P5} are isomorphisms.  They are natural under the
stage transitions, smooth chart changes, and retention of
marked parameters.  The same conclusion holds for the
fixed-point, fixed-label standard Whittaker objects of
\eqref{eq:W3}.
\end{proposition}

\begin{proof}
At exponent zero the two maps are precisely Lin's
geometric comparison
\cite[Propositions~2.2.4--2.2.5, Section~2.2.6,
Section~8.1.3 and Proposition~8.1.4, Theorem~9.0.2]{Lin}.
The restriction map is his vanishing after image from
the nontransverse locus; the first arrow is the diagonal
orientation and transverse $!$-to-$*$ construction of
his transformation.  The theorem says that both are
isomorphisms for ordinary Whittaker arguments, also with
marked parameters retained.  Its diagonal is the
ordinary diagonal; thus its torus-quotient presentation
includes the central-cover push appearing in $J^b$.
Its dimension is the stack dimension $b$, with the
unshifted pair of opposite Whittaker objects used here.

The assertion permits our chart spectator.  To check
this without an assertion about arbitrary tensor
extension, first add extra marked points with zero
pole allowance.  Their pole model is the pullback of
the original one, and the smooth modification tests
preserve the Whittaker condition.  Embed the affine
spectator into affine space.  Smoothly, or \'etale
locally, that space is covered by schemes \'etale over
it which are open in powers of $X$: use coordinate
charts on the curve and translations in affine space.
For such opens the assertion follows from the retained
parameter theorem by open shriek extension to the
whole power of $X$.  Smooth-arrow character tests
preserve equivariance under that extension.  Restrict
back to the open, and then use the closed image of
the spectator.  Proper base change commutes that
closed image with every formula, since the spectator
is carried along.  Closed inclusions of pole bounds
handle changes of bounds in this reduction.

Now express the three terms and two arrows of
\eqref{eq:P5} as finite pre-axis calculations on the
chosen bounded-zero cut, using
$\widehat M_{I,m}$ and the regular constructible
frame-ratio formula for $J^b$.  The latter introduces
a Kummer factor but no exponential.  Multiply its
recording coordinate with that of
$\widehat M_{I,m}$ in \eqref{eq:P2}.  Pullback by $g$
lifts to the frame torsors, and the inner frame
descends because its two recording scalings are
opposite.  Thus all the arrows, including orientation
and the push comparison, have the finite geometric
presentation required by Lemma~\ref{deform:rank}.

Apply $C_!^u$ to the cones of the two arrows.  The
Cartan dilation preserves transversality, so the
Fourier specialization rules identify their field
specializations with the cones in \eqref{eq:P5} on
the punctured Fourier axis.  At $q=1$ these projected
cones vanish by the ordinary-level result and
dilation.  Although the pre-object itself need not
have clean support, after Fourier at every nonzero
character parameter the bounded-zero cut suffices
by the clean support of the Whittaker object.  This
justifies the use of that cut in the test at $q=1$.
Lemma~\ref{deform:rank} makes the projected cones
zero for every non-torsion $q$, in particular for
$q=\exp(2\pi i s)$ with $s\notin\mathbb Q$.
Fourier and restriction at $t=1$ prove the claim.

For a fixed-label standard replace
$\widehat M_{I,m}$ by the formula
\eqref{eq:P3} applied to $\delta_{u=0}$ and its
stratum factor on frame torsors, with its fixed
normalization shift.  The stratum exponentials
extend cleanly as in \eqref{eq:W3}, and the same
finite argument applies.  All maps used were
geometric and compatible with the indicated
base changes, which proves the stated naturality.
\end{proof}

By continuity, the colimit of the stage comparisons
identifies the Fourier-specialized augmentation
\eqref{eq:P4}, for irrational $s$, with a canonical
map
\begin{equation}\label{eq:P6}
                         J_sR_s(F_W)\longrightarrow F_W.
\end{equation}
The projection counits and the Vinberg arrows make
this a map of kernels, compatible with smooth
spectator charts.  At exponent zero the analogous
counit has only an anti-tempered cone.  The next
section applies the finite Whittaker test to that
cone, proves the uniform upper weight bound needed
in Lemma~\ref{deform:colimit}, and thereby proves
that \eqref{eq:P6} is an isomorphism at irrational
exponent.
